\providecommand{\toplevelprefix}{../..}  %
\documentclass[../../book-main.tex]{subfiles}

\begin{document}

\chapter*{2.0 版前言}

% \begin{center}
%     ``{\em 条条大路通罗马}。''

% \end{center}
% \vspace{5mm}

作为一本开源书籍，本书于2025年8月18日以“数据分布的深度表征学习”为题发布的第一版，在过去的六个月里持续频繁的更新。同时，结合我们在2025年秋季香港大学一门新课程中使用本书的经验，以及来自同事、助教和学生的反馈，我们发现了书中许多值得进一步修订和扩充的内容。

因此，我们决定在全新的 \textit{2.0 版} 中，对本书的内容和组织结构进行一些实质性的修改和升级。新版本的题目更新为：\\
\centerline{\em 深度表征学习的原理与实践}\\ 
\centerline{\em \small 或记忆的数学理论}
新版本还使我们能够更明确揭示书中不同章节材料之间紧密的概念和技术联系，因此我们相信相较于第一版，其整体教学价值已得到显著提升。

在我们看来，新版本对这一主题的呈现更加统一、完整且条理清晰。新版本还融入了
最新发展的理论见解，以及越来越多针对真实世界数据的实际应用，这些内容
在其他地方尚未得到妥善记录和系统解释。2.0 版正是对这一现状的及时补充。

\subsection*{2.0 版的主要变化}
\begin{itemize}
    \item 我们将第一版中的第3章拆分为两章，即现在的
        第 \ref{ch:denoising} 章和第 \ref{ch:lossy-compression} 章。新的
        第 \ref{ch:denoising} 章重点关注学习低维分布的去噪过程。我们增加了一个新小节，
        刻画了该过程导致泛化或记忆的条件。新的第 \ref{ch:lossy-compression} 章重点关注
        基于有损编码方法和信息增益最大化原理的表征学习。由于旧版本仅
        说明了如何在监督设置下应用该原理来学习图像表征，新版本增加了一个
        无监督设置的案例研究。它还为流行的无监督学习方法（特别是对比学习和 DINO）
        提供了理论依据。
    \item 在关于通过展开优化设计深度表征的更新后的第 \ref{ch:deep-networks} 章中，我们为白盒 Transformer 架构 CRATE 增加了一个因果版本的原理性推导，这对于处理我们在应用中常见的序列数据（如文本）非常重要。
    \item 在关于一致性表征学习的更新后的第 \ref{ch:consistent-representations} 章中，
        我们增加了一个新小节，进一步阐明了分布学习和表征学习之间的关系。
        这为流行的实用自编码方法（如变分自编码器和表征自编码器）提供了
        理论澄清和依据。
    \item 在更新后的第 \ref{ch:conditional-inference} 章中，我们增加了一个小节，
        从互信息的视角对成对数据的表征学习和条件生成提供了原理性解释。
        它为 CLIP 和交叉注意力（cross attention）等流行的实用方法提供了理论依据。
    \item 扩充后的应用章节（第 \ref{ch:applications} 章）现在包含了
        更多针对真实世界数据（包括自然 2D 图像、自然 3D 物体、人体运动和自然语言）的分布学习和表征学习的详细实现，涵盖了几乎所有流行和实用的设置：监督、弱监督、无监督和条件生成。 
    \item 关于开放方向的最后一章（第 \ref{ch:future} 章）也进行了
        大幅重写。我们试图对不同层次的智能给出一个更清晰的分类，以便我们能够
        阐明（通过本书）我们对智能已经做了什么、理解了什么，以及还有什么尚未理解。
        我们相信，与更高级形式的智能相关的开放性问题可以被很好地提出，从而可以
        通过科学和数学手段进行定性和定量研究，而不是仅仅作为一个神秘的主题或纯粹的哲学话题。
\end{itemize}


\subsection*{2.0 版内容贡献者}
除了作者之外，在本书 {\em 2.0 版} 的准备过程中，许多学生和同事也开始加入这个项目，并为本书的不同部分贡献了宝贵的内容。以下是部分贡献者及其具体贡献的名单（按字母顺序排列）：
\begin{itemize}
    \item Tianzhe Chu：第 \ref{sec:image_classification} 节的实验，以及本书的 AI 助手。
    \item Shenghua Gao 教授：第 \ref{sec:Michelangelo} 节和第 \ref{sec:CUPID} 节。
    \item Bingbing Huang：第 \ref{sec:Michelangelo} 节和第 \ref{sec:CUPID} 节。
    \item Qing Qu 教授：第 \ref{sec:diffusion_memorization} 节。
    \item Shengbang Peter Tong：第 \ref{sec:continuous} 节和第 \ref{sec:RAE-implement} 节。
    \item Chengyu Wang：第 \ref{sec:Michelangelo} 节和第 \ref{sec:CUPID} 节。
    \item Ziyang Robin Wu：第 \ref{sec:CLIP} 节。
    \item Chun-Hsiao Daniel Yeh：第 \ref{sec:auto-encode-img-gen} 节和第 \ref{sec:cond_gen} 节。
    \item Brent Yi：第 \ref{sec:egoallo} 节。
    \item Jingfeng Yang：第 \ref{sec:CLIP} 节。
    \item Yaodong Yu 博士：第 \ref{ch:deep-networks} 章基于其博士论文。
    \item Zibo Zhao 博士：第 \ref{sec:Michelangelo} 节。

\end{itemize}

\end{document}